% =====================================================================
%  5.5 Sammensatte problemer — fasit
% =====================================================================

\fk{Del 1 — Teoretiske problemer}
\fo{1.1} a) sinussetningen (to vinkler og en side) \ b) cosinussetningen (to sider og vinkelen mellom) \ c) cosinussetningen (tre sider) \ d) arealsetningen
\fo{1.2} a) $\sin 60^\circ$, $6 \cdot \frac{\sqrt3}{2} = 3\sqrt3$ \ b) $\sin 45^\circ$, $AC = \frac{3\sqrt3}{\sqrt2/2} = 3\sqrt6$
\fo{1.3} a) $60^\circ$, $45^\circ$ og $30^\circ$ \ b) $\frac{180^\circ - 45^\circ}{2} = 67,5^\circ$ ved både $A$ og $B$. $OA = OB$ (avstand fra sentrum til hjørnene), så trekanten er likebeint.
\fo{1.4} a) Trekant $ABC$, SVS: $AC^2 = 8^2 + 15^2 - 2 \cdot 8 \cdot 15 \cdot \cos 60^\circ = 169$, så $AC = 13$.
b) Trekant $ACD$, SSS: $\cos D = \frac{5^2 + 12^2 - 13^2}{2 \cdot 5 \cdot 12} = 0$, så $D = 90^\circ$.
c) $T = \frac12 \cdot 8 \cdot 15 \cdot \sin 60^\circ + \frac12 \cdot 5 \cdot 12 = 30\sqrt3 + 30$
\fo{1.5} a) Trekant $ABD$ (sinussetningen): $\frac{AD}{\sin 45^\circ} = \frac{6}{\sin 60^\circ}$, så $AD = 2\sqrt6$.
b) Trekant $ADC$: $\angle ADC = 180^\circ - 60^\circ = 120^\circ$. Sinussetningen: $AC = \frac{2\sqrt6 \cdot \sin 120^\circ}{\sin 30^\circ} = 6\sqrt2$.
c) $\angle DAC = 180^\circ - 120^\circ - 30^\circ = 30^\circ = \angle C$, så $DC = AD = 2\sqrt6$. $T = \frac12 \cdot 2\sqrt6 \cdot 2\sqrt6 \cdot \sin 120^\circ = 6\sqrt3$.
\fo{1.6} Sentralvinkelen er $72^\circ$. Trekant $OAB$ er likebeint med toppvinkel $72^\circ$.
a) Halver trekanten: $\sin 36^\circ = \frac{2,0}{R}$, så $R \approx 3,4$ cm.
b) $T = 5 \cdot \frac12 \cdot R^2 \cdot \sin 72^\circ \approx 27,5$ cm$^2$ (også med avrundet $R = 3,4$).
\fo{1.7} Feilen er i overgangen til $DC = 5\sqrt2 \cdot \cos 30^\circ$. Vinkelen $30^\circ$ hører til trekant $ABD$, men $DC$ ligger i trekant $ADC$ der vinkelen ved $C$ er $45^\circ$. Riktig: $DC = 5\sqrt2 \cdot \cos 45^\circ = 5$, og $BC = 5\sqrt3 + 5$.
\fo{1.8} a) Noen ganger. Arealet er $4^2 \cdot \sin v = 16 \sin v$. Med $v = 30^\circ$ er arealet $8$, med $v = 90^\circ$ er det $16$.
b) Noen ganger. Cosinussetningen gir $d_1^2 = 32 - 32\cos v$ og $d_2^2 = 32 + 32\cos v$. De er like bare når $\cos v = 0$, altså i et kvadrat.
c) Aldri. Vinkelen på $60^\circ$ ligger mellom to sider $a$ og $b$. Arealet er $\frac12 ab \sin 60^\circ = \frac{ab}{4}\sqrt3$. Siden $ab$ er et positivt heltall og $\sqrt3$ ikke er en brøk, er arealet aldri et helt tall.
\fo{1.9} a) I rettvinklet trekant $ACD$ er $h = b \sin A$. Sinussetningen gir $b = \frac{c \sin B}{\sin C}$. Da er $h = \frac{c \sin A \sin B}{\sin C}$.
b) $C = 120^\circ$ og $h = \frac{6\sqrt3 \cdot \frac12 \cdot \frac12}{\frac{\sqrt3}{2}} = 3$. Kontroll: trekanten er likebeint, halve grunnlinja er $3\sqrt3$, og $h = 3\sqrt3 \cdot \tan 30^\circ = 3$.
\fo{Sjekk 1} a) cosinussetningen (SSS) \ b) $20^\circ$ \ c) Definisjonen krever en rett vinkel. I en vilkårlig trekant kan vi trekke en høyde og få to rettvinklede trekanter, eller bruke sinus-, cosinus- eller arealsetningen.

\fk{Del 2 — Praktiske problemer}
\fo{2.1} a) Øst ($70^\circ$ er nærmere $90^\circ$ enn $0^\circ$) \ b) $70^\circ + 180^\circ = 250^\circ$ \ c) $150^\circ - 70^\circ = 80^\circ$
\fo{2.2} $h = 45 \cdot \tan 32^\circ \approx 28,1$ m
\fo{2.3} a) $66^\circ$ \ b) $\sin 66^\circ$, $AC = \frac{200 \cdot \sin 64^\circ}{\sin 66^\circ} \approx 196,8$ m \ c) $BC = \frac{200 \cdot \sin 50^\circ}{\sin 66^\circ} \approx 167,7$ m
\fo{2.4} $\angle T = 180^\circ - 58^\circ - 74^\circ = 48^\circ$. Sinussetningen: $AT = \frac{75 \cdot \sin 74^\circ}{\sin 48^\circ} \approx 97,0$ m. Bredden er $AT \cdot \sin 58^\circ \approx 82,3$ m. Samme svar med avrundet $AT$.
\fo{2.5} a) $\angle HPQ = 180^\circ - (160^\circ - 30^\circ) = 50^\circ$ \ b) $HQ^2 = 15^2 + 10^2 - 2 \cdot 15 \cdot 10 \cdot \cos 50^\circ$, så $HQ \approx 11,5$ km \ c) $\cos H = \frac{15^2 + HQ^2 - 10^2}{2 \cdot 15 \cdot HQ}$ gir $H \approx 41,8^\circ$ (samme med avrundet $HQ$) \ d) Kurs fra $H$ til $Q$ er $30^\circ + 41,8^\circ = 71,8^\circ$. Tilbake: $71,8^\circ + 180^\circ = 251,8^\circ$.
\fo{2.6} Feilen er i siste linje. Høyden finnes i den rettvinklede trekanten ved $B$, der vinkelen er $32^\circ$, ikke $20^\circ$. Riktig: $h = 493,5 \cdot \sin 32^\circ \approx 261,5$ m. Kontroll: $\frac{h}{\tan 20^\circ} - \frac{h}{\tan 32^\circ} = 300$.
\fo{2.7} Toppvinkelen er $95^\circ$. Venstre sperre (overfor $50^\circ$): $\frac{10 \cdot \sin 50^\circ}{\sin 95^\circ} \approx 7,7$ m. Høyre sperre (overfor $35^\circ$): $\frac{10 \cdot \sin 35^\circ}{\sin 95^\circ} \approx 5,8$ m. Mønehøyde: $7,7 \cdot \sin 35^\circ \approx 4,4$ m.
\fo{2.8} a) $80 \cdot \tan 25^\circ \approx 37,3$ m. Fra $80 \cdot \tan 24^\circ \approx 35,6$ m til $80 \cdot \tan 26^\circ \approx 39,0$ m. Spenn: ca. $3,4$ m.
b) $80 \cdot \tan 60^\circ \approx 138,6$ m. Fra $133,1$ m til $144,3$ m. Spenn: ca. $11,2$ m.
c) Ved $60^\circ$ er siktelinja både lengre og brattere. Når vinkelen endres med $1^\circ$, flytter siktelinja seg mer over masten jo lenger den er. Da endres høyden mer.
\fo{2.9} a) $90^\circ - 30^\circ = 60^\circ$.
b) Skip $A$ har seilt $60$ km. La $x$ være distansen til $B$. Cosinussetningen: $52^2 = 60^2 + x^2 - 2 \cdot 60 \cdot x \cdot \cos 60^\circ$, altså $x^2 - 60x + 896 = 0$. Da er $x = 28$ eller $x = 32$. Farten er $14$ km/t eller $16$ km/t. Begge er mulige.
c) Dette er SSV: to sider ($60$ og $52$) og en vinkel som ikke ligger mellom dem. Begge løsningene for $x$ er positive, og begge gir en trekant. Det er to trekanter, som i svingarmen fra 5.3.
\fo{Sjekk 2} a) $130^\circ$ \ b) $30 \cdot \tan 45^\circ = 30$ m \ c) Avrunding underveis gir avrundingsfeil som vokser for hvert steg. Lagrede verdier er mer nøyaktige.

\fk{Ekstra utfordring}
\fo{U1} a) Åtte trekanter med to sider $1$ og vinkelen $45^\circ$: $T = 8 \cdot \frac12 \cdot 1 \cdot 1 \cdot \sin 45^\circ = 2\sqrt2$.
b) Cosinussetningen: $s^2 = 1 + 1 - 2\cos 45^\circ = 2 - \sqrt2$, så $s = \sqrt{2 - \sqrt2}$.
c) Omkrets: $8\sqrt{2 - \sqrt2} \approx 6,12$. Sirkelen har omkrets $2\pi \approx 6,28$. Hver side er en rett linje, og den er kortere enn buen mellom de samme to punktene.
\fo{U2} a) $50 \cdot \tan 39^\circ \approx 40$ m og $50 \cdot \tan 41^\circ \approx 43$ m (midt på: $42$ m).
b) To gjeldende siffer, altså ca. $42$ m. Spennet er ca. $3$ m, så flere siffer ville gitt en falsk nøyaktighet.
c) Minst: $49,5 \cdot \tan 39^\circ \approx 40$ m. Størst: $50,5 \cdot \tan 41^\circ \approx 44$ m.
\fo{U3} a) $5^2 = (4\sqrt2)^2 + c^2 - 2 \cdot 4\sqrt2 \cdot c \cdot \cos 45^\circ$ gir $25 = 32 + c^2 - 8c$, altså $c^2 - 8c + 7 = 0$. Da er $c = 1$ eller $c = 7$.
b) Begge løsningene er positive, så begge gir en trekant. Siden $a = 5$ er mindre enn $b$ men større enn $b \sin A = 4$, treffer svingarmen grunnlinja to ganger.
c) $\sin B = \frac{4\sqrt2 \cdot \sin 45^\circ}{5} = \frac45$, så $B \approx 53,1^\circ$ eller $B \approx 126,9^\circ$. Da er $C \approx 81,9^\circ$ eller $C \approx 8,1^\circ$, og $c = \frac{5 \sin C}{\sin 45^\circ}$ gir $7,0$ og $1,0$.
\fo{U4} a) Betingelsen er \lstinline[style=laereverk,numbers=none,frame=none]|A + B2 < 180|. Utskrift: \lstinline[style=laereverk,numbers=none,frame=none]|60.0|, \lstinline[style=laereverk,numbers=none,frame=none]|7.0|, \lstinline[style=laereverk,numbers=none,frame=none]|[53.1, 126.9]|, \lstinline[style=laereverk,numbers=none,frame=none]|[26.4]| og en tom liste.
b) $\sin B = \sin(180^\circ - B)$, så det finnes to vinkler som passer. Den andre gir bare en trekant hvis vinkelsummen er under $180^\circ$. Tom liste betyr $\sin B > 1$: ingen trekant.
c) \lstinline[style=laereverk,numbers=none,frame=none]|def trekant_sss(a, b, c):| gir \lstinline[style=laereverk,numbers=none,frame=none]|A = vinkel_fra_sider(a, b, c)|, \lstinline[style=laereverk,numbers=none,frame=none]|B = vinkel_fra_sider(b, a, c)|, \lstinline[style=laereverk,numbers=none,frame=none]|C = 180 - A - B| og \lstinline[style=laereverk,numbers=none,frame=none]|return [A, B, C]|.
\fo{U5} a) $1 + \cos C = \frac{(a+b)^2 - c^2}{2ab}$ og $(a+b)^2 - c^2 = (a+b-c)(a+b+c)$.
b) Faktorene er $a+b+c = 2s$, $a+b-c = 2(s-c)$, $c+a-b = 2(s-b)$ og $c-a+b = 2(s-a)$. Produktet er $16\,s(s-a)(s-b)(s-c)$. Delt på $16$ gir $T^2 = s(s-a)(s-b)(s-c)$. Ta kvadratroten.
c) $s = 10$ og $T = \sqrt{10 \cdot 5 \cdot 3 \cdot 2} = \sqrt{300} = 10\sqrt3$. Kontroll: vinkelen overfor $7$ har $\cos = \frac{25 + 64 - 49}{80} = \frac12$, altså $60^\circ$. $T = \frac12 \cdot 5 \cdot 8 \cdot \sin 60^\circ = 10\sqrt3$.
\fo{U6} Flere svar er riktige. Eksempel: firkant $ABCD$ med $AB = 7$, $BC = 8$, $\angle B = 120^\circ$, $\angle ACD = 90^\circ$ og $CD = 84$. Finn $AD$. Cosinussetningen: $AC^2 = 49 + 64 + 56 = 169$, så $AC = 13$. Pytagoras: $AD = \sqrt{13^2 + 84^2} = 85$.
a) Vinkler med eksakt cosinus som er en brøk, som $60^\circ$ og $120^\circ$ ($\cos = \pm\frac12$).
\fo{U7} a) $T = \frac12 \cdot x \cdot (20 - x) \cdot \sin 30^\circ = \frac14 x(20 - x)$.
b) $T(x) = -\frac14 x^2 + 5x$. Toppunktet er ved $x = \frac{-5}{2 \cdot (-\frac14)} = 10$. Derivasjon: $T'(x) = -\frac12 x + 5 = 0$ gir $x = 10$. Størst areal: $T(10) = 25$ m$^2$.
c) $BC^2 = 10^2 + 10^2 - 2 \cdot 10 \cdot 10 \cdot \cos 30^\circ = 200 - 100\sqrt3$. $BC = 10\sqrt{2 - \sqrt3} \approx 5,2$ m (også lik $5(\sqrt6 - \sqrt2)$).
